% Provenance (never printed): see sources/outline_sources.md (Year 3:
% stereographic projection; fold and fault analysis; palaeostress; balanced
% cross-sections). Brief: sources/book5/BRIEFS.md, ch. 8. Nets, pi-diagram and
% fault-slip data from figdata/bachelor-3/stereonet.py (Varden data invented,
% from figdata/shared/region.py); Fry plot from figdata/bachelor-3/fry.py.
% Exercise and problem values in the Book 5 scratch calculator calc/ch08.py.
% No real-world numbers: no ledger rows.
\chapter{构造分析}\label{ch:b3:structural-analysis}

一位地质学家在一个大到无法整体看清的褶皱两翼上
测得两百个层面：转折端埋在森林之下，两翼
延伸数千米。每个读数都是一对数。将它们画成圆形图上的点，
层面的极沿一条单一曲线排列，该曲线即一个大圆，而
该大圆的极即褶皱轴，这是一条
从未有人见过的直线。本章建立这张图，即\omterm{def:b3:structural-analysis:stereo}{赤平投影}，
并借助其背后的矢量工具，求褶皱轴、测量
变形物体所记录的应变，恢复驱动一组断层的
应力，并将一个变形\omterm{def:b3:structural-analysis:restoration}{剖面恢复}到变形之前的
状态。

\begin{recall}
球面上的大圆与小圆见本科一年级分册中
关于球面天文的一章；走向、倾角、倾向、倾伏向与倾伏角见其
关于地质图与剖面的章节；倾向为
$\alpha$、倾角为 $\delta$ 的平面，其向上单位法向量为 $\vec n = (\sin\delta
\sin\alpha, \sin\delta\cos\alpha, \cos\delta)$（东、北、上分量），
与一组读数的矢量平均，见其第一章的
野外与实验方法；褶皱要素与类型，
褶皱的枢纽平行于 $\vec n_1\times\vec n_2$，两翼间夹角为
$\ang{180} - \arccos(\vec n_1\cdot\vec n_2)$，以及安德森的
断层理论见其断层、褶皱与组构一章；应力、
主应力、应变椭圆与简单剪切，见其
应力、应变与流变学一章；滑动矢量的侧伏角（自走向起算）
正断层取 \ang{-90}，逆断层取 \ang{90}，左行取 0，
右行取 \ang{180}）；震源机制的 P 轴与 T 轴，见其
地震源一章。平衡剖面与滑脱面见
本科二年级分册中关于阿尔卑斯山的一章。最小二乘法见
\cref{sec:b3:seismology-imaging:inverse}；对称矩阵的本征值与特征向量
见数学类书籍。
\end{recall}

\section{赤平投影}\label{sec:b3:structural-analysis:stereo}

地质学家的数据是方位：层面的产状、
断层的产状、面理的产状；线理、褶皱转折端与滑动方向的倾伏向和倾伏角。
每一项都是过某点的方向，本身并无位置。为比较数以百计的方向，
设想把它们全部平移到同一球面的球心：
一条直线与球的下半球相交于一点，一个平面则沿
一条半大圆与之相交。于是下半球的图上，直线表现为
点，平面表现为曲线。

\begin{definition}[赤平投影，等角网，等面积投影，极]\label{def:b3:structural-analysis:stereo}
设过球心 $O$、半径为 $R$ 的球的一条直线，与下半球
相交于 $P$，其与天底的夹角为 $\theta = \ang{90} - p$，其中
$p$ 为该直线的倾伏角，$t$ 为其倾伏向。\emph{赤平投影}\index{赤平投影}
将 $P$ 映为由天顶至 $P$ 的直线
与水平面的交点，该水平面过
$O$；该点距离为 $r$，自 $O$ 量起，方位角为 $t$。
按此方式画出的下半球图，连同其由大圆与小圆组成的网格，
称为\emph{等角网}\index{等角网}（或
Wulff 网）。另有球面的\emph{等面积投影}\index{等面积投影}，
将 $P$ 映为距离为 $r =
R\sqrt2\sin(\theta/2)$、方位角为 $t$ 的点；其网为施密特网。
平面的\emph{极}\index{极}是其法线，画作
一个点：倾向为 $\alpha$、倾角为 $\delta$ 的平面，其极
的倾伏向为 $\alpha + \ang{180}$，倾伏角为 $\ang{90} - \delta$。
\end{definition}

两种投影都把竖直线映到中心，把水平线映到
半径为 $R$ 的圆（即基圆）上。二者的区别在于保留的性质不同。

\begin{theorem}[两种投影]\label{thm:b3:structural-analysis:projections}
\begin{enumerate}[label=(\alph*)]
  \item \omterm{def:b3:structural-analysis:stereo}{赤平投影}将 $P$ 置于 $r = R\tan(\theta/2)$ 处。它
        是保角的：在每一点，它各方向的伸缩倍数相同，
        均为 $1/(2\cos^2(\theta/2))$，因而保持
        曲线之间的夹角。
  \item \omterm{def:b3:structural-analysis:stereo}{等面积投影}使每个面积都乘以同一个因子 $1/2$：
        网上某一区域的面积，是其所代表的
        半球区域面积的一半。
\end{enumerate}
\end{theorem}

\begin{proof}
(a) 在过天顶 $Z$、天底与 $P$ 的竖直平面内
（见下方第一幅图），由天底至 $P$ 的弧
在球心处所张的角为 $\theta$，因而圆周角为
$\theta/2$，顶点为 $Z$，其高度为 $R$，在 $O$ 之上：$r =
R\tan(\theta/2)$。球面上朝天底方向的一小步 $R\,\mathrm d\theta$
在图上变为 $\mathrm dr = R\,\mathrm d\theta/(2\cos^2(\theta/2))$。
与天底成 $\theta$ 角的各点构成的圆，其周长为
$2\pi R\sin\theta$，在图上变为周长为 $2\pi
r$ 的圆；二者之比为 $\tan(\theta/2)/\sin\theta = 1/(2\cos^2(\theta/2))$，
与沿半径方向的比值相同。若映射在两个互相垂直的方向上伸缩相同，
则它在一切方向上伸缩都相同，因而保持角度。

(b) 球面上与天底夹角在 $\theta$ 之内的球冠，其面积为
$2\pi R^2(1 - \cos\theta) = 4\pi R^2\sin^2(\theta/2)$；其像为
面积为 $\pi r^2 = 2\pi R^2\sin^2(\theta/2)$ 的圆盘，恰为其一半。两者的面积
在方位角上均匀分配，故两方位角之间的球冠扇形与其像
保持比值 $1/2$；两个值之间的环带亦然，
即两个 $\theta$ 所对球冠之差。每个区域都是若干此类
小环扇形之并。
\end{proof}

\omterm{def:b3:structural-analysis:stereo}{赤平投影}还将球面上的每个圆，无论大圆还是小圆，
映为图上的一个圆或一条直线；此事在本阶段不加证明地承认。它使
\omterm{def:b3:structural-analysis:stereo}{等角网}便于用圆规画出；小圆，即与某轴成固定夹角的直线所构成的圆锥，
在图上保持其形状。\omterm{def:b3:structural-analysis:stereo}{等面积投影}则
会扭曲形状：小圆变为一个椭圆，沿
基圆伸展。它是可以对一簇极进行计数的投影，
因为网上面积相等的区域包含半球方向上相等的份额。
构造地质学家与为震源机制工作的地震学家，
都用它来处理各自的数据。

\begin{omfigure}
\begin{tikzpicture}[scale=1.0, >=stealth]
% (a) equal-angle: theta = 60 deg, R = 2 cm; r = 2 tan 30 = 1.155
\begin{scope}
\draw[thick] (0,0) circle (2);
\draw (-2.5,0) -- (2.5,0) node[right, font=\scriptsize] {水平面};
\fill (0,2) circle (1.5pt) node[above, font=\scriptsize] {天顶 $Z$};
\fill (0,-2) circle (1.5pt) node[below, font=\scriptsize] {天底};
\fill (0,0) circle (1.2pt) node[above left, font=\scriptsize] {$O$};
\draw[dashed] (0,2) -- (0,-2);
\draw[omDef, thick] (0,0) -- (1.732,-1);
\fill[omDef] (1.732,-1) circle (1.5pt) node[below right, font=\scriptsize] {$P$};
\draw[omThm] (0,2) -- (1.732,-1);
\fill[omThm] (1.155,0) circle (1.5pt);
\node[omThm, font=\scriptsize, above right] at (1.155,0) {$P'$};
\draw (0,-0.6) arc[start angle=-90, end angle=-30, radius=0.6];
\node[font=\scriptsize] at (0.38,-0.85) {$\theta$};
\draw (0,1.2) arc[start angle=-90, end angle=-60, radius=0.8];
\node[font=\scriptsize] at (0.25,0.95) {$\theta/2$};
\draw[<->] (0,0.25) -- (1.155,0.25);
\node[font=\scriptsize, above] at (0.58,0.25) {$r$};
\node[font=\small] at (0,-3.0) {(a) 等角投影};
\end{scope}
% (b) equal-area: chord from nadir to P = 2R sin(theta/2) = 2
\begin{scope}[xshift=6.6cm]
\draw[thick] (0,0) circle (2);
\draw (-2.5,0) -- (2.5,0);
\draw (-2.5,-2) -- (2.5,-2) node[right, font=\scriptsize] {切面};
\fill (0,-2) circle (1.5pt) node[below left, font=\scriptsize] {天底};
\fill (0,0) circle (1.2pt) node[above left, font=\scriptsize] {$O$};
\draw[dashed] (0,2) -- (0,-2);
\draw[omDef, thick] (0,0) -- (1.732,-1);
\fill[omDef] (1.732,-1) circle (1.5pt) node[above right, font=\scriptsize] {$P$};
\draw[omThm] (0,-2) -- (1.732,-1);
\draw[omThm, dashed] (1.732,-1) arc[start angle=30, end angle=0, radius=2];
\fill[omThm] (2,-2) circle (1.5pt);
\node[omThm, font=\scriptsize, below] at (2,-2.05) {$2R\sin(\theta/2)$};
\draw (0,-0.6) arc[start angle=-90, end angle=-30, radius=0.6];
\node[font=\scriptsize] at (0.38,-0.85) {$\theta$};
\node[font=\small] at (0,-3.0) {(b) 等面积投影};
\end{scope}
\end{tikzpicture}

\par\smallskip{\small 两种投影都画在过直线
$OP$ 的竖直平面内。（a）\omterm{def:b3:structural-analysis:stereo}{赤平投影}将 $P$ 映为 $P'$，即
自天顶引出的直线与水平面的交点：在 $Z$ 处的角为 $\theta/2$，故
$r = R\tan(\theta/2)$。（b）\omterm{def:b3:structural-analysis:stereo}{等面积投影}绕天底将
$P$ 投到切面上，该切面位于弦长 $2R\sin(\theta/2)$ 处；
图随后按 $1/\sqrt2$ 缩放，使水平线
落在半径为 $R$ 的基圆上。}\label{fig:b3:structural-analysis:geometry}
\end{omfigure}

\begin{omfigure}
\begin{tikzpicture}
\begin{groupplot}[group style={group size=2 by 1, horizontal sep=0.8cm},
  width=0.5\linewidth, height=0.5\linewidth, axis equal image, hide axis,
  xmin=-1.15, xmax=1.15, ymin=-1.15, ymax=1.15, unbounded coords=jump,
  title style={font=\small}]
\nextgroupplot[title={(a) 等角网（Wulff）}]
% computed in figdata/bachelor-3/stereonet.py
\addplot[gray!75, very thin] table[x=x, y=y] {figdata/out/bachelor-3-stereonet-wulff.dat};
\addplot[gray!75, very thin] coordinates {(-1,0) (1,0)};
\addplot[black, thick, domain=0:360, samples=181] ({cos(x)}, {sin(x)});
\node[font=\scriptsize, above] at (axis cs:0,1.0) {北};
\nextgroupplot[title={(b) 等面积网（Schmidt）}]
\addplot[gray!75, very thin] table[x=x, y=y] {figdata/out/bachelor-3-stereonet-schmidt.dat};
\addplot[gray!75, very thin] coordinates {(-1,0) (1,0)};
\addplot[black, thick, domain=0:360, samples=181] ({cos(x)}, {sin(x)});
\node[font=\scriptsize, above] at (axis cs:0,1.0) {北};
\end{groupplot}
\end{tikzpicture}

\par\smallskip{\small 两张网依据上文定理的公式绘制：大圆对应走向南北、
分别向东和向西每隔
\ang{10} 倾斜的平面，另有绕南北轴、间隔为 \ang{10} 的小圆。在\omterm{def:b3:structural-analysis:stereo}{等角网}上，
小圆仍为圆；在等面积网上，小圆变为椭圆，
但等面积区域所含的
方向份额相等。}\label{fig:b3:structural-analysis:nets}
\end{omfigure}

\paragraph{平面、直线与夹角。}
倾向为 $\alpha$、倾角为 $\delta$ 的平面在下半球上截出一条
半大圆。其两端位于基圆上，即走向的两端；
其中点为最大倾斜线，倾伏角为 $\delta$，
朝 $\alpha$ 方向倾伏；其极点与其上每一点相距 \ang{90}，位于
方位角 $\alpha + \ang{180}$ 的方向上。平面内的直线在网上
对应大圆上的一点。作图时，将透明覆盖纸铺在网上，
绕圆心转动，直至数据的走向沿网的南北
直径方向；随后描出倾角相符的大圆，再将覆盖纸
转回。两直线之间的夹角沿过两点的大圆量取，
数其所跨越的小圆分度，此前须
转动覆盖纸，使两点位于网上同一大圆上。
两平面之间的夹角即其极点间的夹角，或其补角。
计算机以矢量完成同样的计算：两直线的单位矢量
$\vec u$ 与 $\vec v$ 之间的夹角为 $\arccos|\vec u\cdot\vec v|$。

\begin{method}[绕轴旋转数据]\label{met:b3:structural-analysis:rotation}
为消除倾斜，例如使岩层连同其所含的
线理恢复至水平，或将倾斜地层中测得的古地磁方向
复原（见\cref{ch:b3:palaeomagnetism}）：
\begin{enumerate}
  \item 选定旋转轴：对于倾斜的岩层，取其走向线，
        水平；旋转角即倾角，旋转方向使岩层抬起。
  \item 在网上转动覆盖纸，使旋转轴沿
        南北向直径，此时小圆即绕它的
        圆锥面。
  \item 使每一点沿其小圆移动旋转角，
        旋转角沿大圆量取。若某点将离开下半球，
        则它在该小圆的另一端重新出现，
        且位于相反的方位角上。
  \item 将覆盖纸转回。用矢量表示时，$\vec v$ 转过
        角 $\omega$，绕单位轴 $\vec k$ 的旋转为 $\vec v' = \vec v
        \cos\omega + (\vec k\times\vec v)\sin\omega + \vec k(\vec k\cdot\vec
        v)(1 - \cos\omega)$。
\end{enumerate}
\end{method}

\section{褶皱分析}\label{sec:b3:structural-analysis:folds}

圆柱状褶皱由一条直线（即褶皱轴）
沿自身平行移动生成。与褶皱面相切的每个平面都包含
该轴，因此褶皱面的每条法线都垂直于该轴：所有
极点都落在与该轴垂直的平面所对应的大圆上。

\begin{definition}[π 图解、β 轴]\label{def:b3:structural-analysis:pi}
褶皱的 \emph{π 图解}\index{π 图解}是对其褶皱面上所测平面的
极点作的图。对圆柱状褶皱，这些极点落在
一个大圆上，即 $\pi$ 环带；环带的极点即
\emph{β 轴}\index{β 轴}，是由测量给出的
褶皱轴估计。
\end{definition}

只有两翼时，\omterm{def:b3:structural-analysis:pi}{β 轴}即两翼法线的叉积，
正如本科一年级分册对枢纽所作的说明。若有两百个散布的读数，
则需要最佳拟合。矢量平均曾用于单个平面的读数，
但不再适用：环带上的极点沿其各方向排列，并且
它们的平均几乎说明不了所在的大圆。

\begin{definition}[方位张量]\label{def:b3:structural-analysis:orientation-tensor}
\emph{方位张量}\index{方位张量}是由 $N$ 个单位矢量
$\vec p_1, \dots, \vec p_N$ 给出的对称矩阵为
\[
  T = \frac1N\sum_{i=1}^N \vec p_i\,\vec p_i^{\mathsf T},
  \qquad T_{jk} = \frac1N\sum_i p_{ij}\,p_{ik} .
\]
将矢量换成其相反矢量时，它不变，因此适用于
轴与极点，二者都没有方向。其本征值 $\lambda_1 \ge \lambda_2
\ge \lambda_3 \ge 0$ 之和等于其迹，即 $1$。
\end{definition}

\begin{theorem}[由方位张量得到的聚类与环带]\label{thm:b3:structural-analysis:tensor}
对任一单位矢量 $\vec u$，
\[
  \frac1N\sum_{i=1}^N (\vec p_i\cdot\vec u)^2 = \vec u^{\mathsf T}T\,\vec u ,
\]
该式取最小值 $\lambda_3$ 时，$\vec u$ 为
最小本征值对应的特征向量 $\vec e_3$，其最大值 $\lambda_1$ 则沿
特征向量 $\vec e_1$ 取得。由于 $\vec p_i\cdot\vec u$ 等于
$\vec p_i$ 与垂直于 $\vec u$ 的平面间夹角的正弦，故
与 $\vec e_3$ 垂直的大圆，即使各极点到它的
距离的正弦平方和最小的环带，而 $\sqrt{\lambda_3}$ 即其
正弦的方均根。轴 $\vec e_1$ 是与这些矢量最接近的直线，
即聚类的平均轴。
\end{theorem}

\begin{proof}
$(\vec p_i\cdot\vec u)^2 = \vec u^{\mathsf T}\vec p_i\,\vec p_i^{\mathsf T}\vec
u$；对 $i$ 取平均。矩阵 $T$ 是对称的，故有一组标准正交的
特征向量基，在该基下 $\vec u = \sum_k u_k\vec e_k$，其中 $\sum_k
u_k^2 = 1$，且 $\vec u^{\mathsf T}T\vec u = \sum_k\lambda_k u_k^2$。这是
特征值的加权平均，权重 $u_k^2$ 之和为 1，因此它
介于 $\lambda_3$ 与 $\lambda_1$ 之间，并在 $\vec u = \vec
e_3$ 与 $\vec e_1$ 时取到这两个值。
\end{proof}

本征值刻画了数据的形状。聚类的 $\lambda_1$ 接近
1，另外两个较小；环带的 $\lambda_3$ 接近 0，极点
沿该圆分布，使 $\lambda_1$ 与 $\lambda_2$ 分担其余；
均匀散布的矢量给出三个接近 $1/3$ 的本征值。在读数精度范围内，褶皱
可视为圆柱状，当 $\sqrt{\lambda_3}$ 不大于
同一平面重复读数的离散度。圆锥状褶皱，
或两期褶皱，会使极点偏离任一大圆，
并使 $\lambda_3$ 增大。

下图绘出本系列虚构的瓦尔登
背斜两翼的十二个读数（见\cref{pb:b3:structural-analysis:1}，
可求出其褶皱轴）。缓倾两翼的极点聚集在网心附近，
分为两组，二者共同确定了一个几乎呈东西向的环带；
\omterm{def:b3:structural-analysis:pi}{β 轴}位于基圆上，几乎正北。两翼的平均
平面以大圆画出，三者相交于同一点：
\omterm{def:b3:structural-analysis:pi}{β 轴}即为所有翼共有的直线。

\begin{omfigure}
\begin{tikzpicture}
\begin{groupplot}[group style={group size=2 by 1, horizontal sep=1.0cm},
  width=0.5\linewidth, height=0.5\linewidth, axis equal image,
  unbounded coords=jump, title style={font=\small}]
\nextgroupplot[hide axis, xmin=-1.15, xmax=1.15, ymin=-1.15, ymax=1.15,
  title={(a) 下半球}]
% computed in figdata/bachelor-3/stereonet.py (Varden readings invented)
\addplot[gray!35, very thin] table[x=x, y=y] {figdata/out/bachelor-3-stereonet-schmidt.dat};
\addplot[black, thick, domain=0:360, samples=181] ({cos(x)}, {sin(x)});
\addplot[gray, densely dashed] table[x=x, y=y] {figdata/out/bachelor-3-stereonet-limbs.dat};
\addplot[omDef, thick] table[x=x, y=y] {figdata/out/bachelor-3-stereonet-girdle.dat};
\addplot[only marks, mark=*, mark size=1.2pt] table[x=x, y=y] {figdata/out/bachelor-3-stereonet-poles.dat};
\addplot[only marks, mark=triangle*, mark size=3pt, omThm] table[x=x, y=y] {figdata/out/bachelor-3-stereonet-beta.dat};
\node[font=\scriptsize, above] at (axis cs:0,1.02) {北};
\node[omThm, font=\scriptsize, right] at (axis cs:0.06,0.88) {$\beta$};
\draw[gray] (axis cs:-0.3,-0.3) rectangle (axis cs:0.3,0.3);
\nextgroupplot[xmin=-0.3, xmax=0.3, ymin=-0.3, ymax=0.3, title={(b) 中心，放大},
  xtick=\empty, ytick=\empty]
\addplot[gray!35, very thin] table[x=x, y=y] {figdata/out/bachelor-3-stereonet-schmidt.dat};
\addplot[omDef, thick] table[x=x, y=y] {figdata/out/bachelor-3-stereonet-girdle.dat};
\addplot[only marks, mark=*, mark size=1.5pt] table[x=x, y=y] {figdata/out/bachelor-3-stereonet-poles.dat};
\node[font=\scriptsize, align=center] at (axis cs:0.16,0.09) {西翼的 \\极点};
\node[font=\scriptsize, align=center] at (axis cs:-0.16,0.09) {东翼的 \\极点};
\node[omDef, font=\scriptsize] at (axis cs:0.17,-0.06) {$\pi$ 环带};
\end{groupplot}
\end{tikzpicture}

\par\smallskip{\small 瓦尔登背斜的 \omterm{def:b3:structural-analysis:pi}{π 图解}，绘于等面积
网（下半球）：十二个层理读数的极点（圆点）、
由\omterm{def:b3:structural-analysis:orientation-tensor}{方位张量}得到的最佳拟合环带（蓝色），两翼的
平均大圆（虚线）与 \omterm{def:b3:structural-analysis:pi}{β 轴}（三角形），三条大圆
相交于同一点。向东倾的岩层，其极点位于
网心以西。}\label{fig:b3:structural-analysis:pi}
\end{omfigure}

\section{有限应变分析}\label{sec:b3:structural-analysis:strain}

一块变形岩石若含有原始形状或排列已知的物体，
便能记录其应变：圆形鲕粒、还原斑点、
对称性已知的化石、卵石，
或原本等间距排列的颗粒中心。

\begin{definition}[有限应变分析，弗赖伊法]\label{def:b3:structural-analysis:strain-analysis}
\emph{有限应变分析}\index{有限应变分析}是测定岩石所经历的应变椭圆（或应变椭球），
依据的是变形前状态已知的物体的形状、方位或间距。
在\emph{弗赖伊法}\index{弗赖伊法}中，
自同一原点画出每个物体中心指向其他各中心的矢量。
若各中心原本以最小间距分布，且无优选方向，
则原点周围的空白区域
即为与应变椭圆形状和方向
相同的椭圆。
\end{definition}

\begin{proposition}[弗赖伊空区]\label{prop:b3:structural-analysis:fry}
若齐次应变将每个矢量 $\vec d$ 映为 $F\vec d$，且应变前任意两中心相距不小于 $d_0$，
则应变后，弗赖伊图中原点附近无点分布的区域
包含椭圆
$\{F\vec d : |\vec d| < d_0\}$，其半轴分别为 $S_1d_0$ 与 $S_2d_0$，沿应变的
主方向。
\end{proposition}

\begin{proof}
齐次应变对每个点施以相同的线性映射再加平移，
故两中心之差变为 $F\vec d$。应变前，没有差值落在半径为 $d_0$ 的圆盘内，
应变后，也就没有差值落在其像之内；
该像即本科一年级分册中的应变椭圆按 $d_0$ 放大所得。
\end{proof}

实际中，各中心并不都处于最小间距，空白区域应读作最密一圈点的内缘，
据此估得的比值略偏低，即应变比 $R_s = S_1/S_2$ 的估计。下图以一张虚构的薄片作检验，
该薄片应变至 $R_s = 2$，长轴与水平线成 \ang{30}：
拟合于空白区域内缘的椭圆给出
$R_s = 1.9$ 与 \ang{30}。
若直接测量物体本身，原为圆形的物体
恰好变为轴比为 $R_s$ 的椭圆；原为椭圆形的物体，
其方向随机，末态比值 $R_f$ 散布于
$R_s$ 附近，原长轴与拉伸方向平行者比值最大，
末态方向 $\varphi$ 则聚集于拉伸方向附近。
以 $R_f$ 对 $\varphi$ 作图，借助表格或计算机读取，
即可将应变与原始形状分开。

\begin{omfigure}
\begin{tikzpicture}
\begin{groupplot}[group style={group size=2 by 1, horizontal sep=1.2cm},
  width=0.5\linewidth, height=0.5\linewidth, axis equal image,
  title style={font=\small}, tick label style={font=\scriptsize}]
\nextgroupplot[xmin=-10, xmax=10, ymin=-10, ymax=10, title={(a) 变形后的颗粒中心}]
% computed in figdata/bachelor-3/fry.py (invented population)
\addplot[only marks, mark=o, mark size=1.6pt, black] table[x=x, y=y] {figdata/out/bachelor-3-fry-points.dat};
\nextgroupplot[xmin=-2.6, xmax=2.6, ymin=-2.6, ymax=2.6, title={(b) 弗赖伊图}]
\addplot[only marks, mark=*, mark size=0.5pt, black] table[x=x, y=y] {figdata/out/bachelor-3-fry-fry.dat};
\addplot[omThm, thick] table[x=x, y=y] {figdata/out/bachelor-3-fry-fit.dat};
\end{groupplot}
\end{tikzpicture}

\par\smallskip{\small 对虚构颗粒中心群体的\omterm{def:b3:structural-analysis:strain-analysis}{弗赖伊法}：颗粒中心先以最小间距 1 排布，
再应变至 $R_s = 2$，长轴自水平线逆时针转 \ang{30}，面积保持不变。（a）
颗粒中心，肉眼几乎看不出应变。
（b）两两距离小于 2.6 的所有中心对之间的矢量，
中心处的空白区域即为应变椭圆。图中红色椭圆
拟合于 36 个扇区中各自最近的点，
比值为 1.9，长轴位于 \ang{30}。}\label{fig:b3:structural-analysis:fry}
\end{omfigure}

三维情形下，应变椭球的主伸长为 $X \ge Y
\ge Z$。其形状绘于 弗林图，以 $a = X/Y$ 对 $b =
Y/Z$ 作图，并引入参数
\[
  k = \frac{a - 1}{b - 1} .
\]
在直线 $k = 1$ 上，即体积不变的平面应变（$Y = 1$，
$X = 1/Z$，故 $a = b$），椭球既非雪茄形，也非煎饼形。直线以上，
$k > 1$，椭球为长球形，因收缩而沿一个方向拉长：
以此方式变形的岩石形成显著的线理。直线以下，$k < 1$，
椭球为扁球形，被压扁，岩石形成显著的面理。

\begin{omfigure}
\begin{tikzpicture}[scale=1.0, >=stealth]
\draw[->] (0,0) -- (5.4,0) node[right, font=\scriptsize] {$b = Y/Z$};
\draw[->] (0,0) -- (0,4.6) node[above, font=\scriptsize] {$a = X/Y$};
\foreach \v/\l in {0/1, 1.5/2, 3/3, 4.5/4} {
  \draw (\v,0) -- (\v,-0.08) node[below, font=\scriptsize] {\l};
  \draw (0,\v) -- (-0.08,\v) node[left, font=\scriptsize] {\l};
}
\draw[omDef, thick] (0,0) -- (4.5,4.5) node[right, font=\scriptsize, align=left] {$k = 1$\\ 平面应变};
\node[font=\scriptsize, align=center] at (1.3,3.6) {收缩 \\ $k > 1$，长球形};
\node[font=\scriptsize, align=center] at (4.3,1.1) {压扁 \\ $k < 1$，扁球形};
% shapes
\draw[fill=omLithLime] (0.9,2.75) ellipse (0.55 and 0.12);
\draw[fill=omLithLime] (4.3,2.1) ellipse (0.45 and 0.22);
\end{tikzpicture}

\par\smallskip{\small 应变椭球的 弗林图。未变形状态位于原点，
即 $a = b = 1$；收缩形成的雪茄形椭球位于
靠近竖轴的一侧，压扁形成的煎饼形椭球位于
靠近横轴的一侧。}\label{fig:b3:structural-analysis:flinn}
\end{omfigure}

\section{古应力}\label{sec:b3:structural-analysis:palaeostress}

安德森理论指出新断层在哪些面上形成，
并指出它们沿垂直于 $\sigma_2$ 的方向滑动。不过，多数断层是较早形成的
薄弱面，在并非其自身所选的应力场中重新活动，
因而常呈斜向滑动。读取断层面上擦痕所示的滑动方向，
仍能记录应力，
前提是每条断层都沿应力推动它的方向滑动。

\begin{definition}[古应力，华莱士–博特假设，应力比，断层滑动反演]\label{def:b3:structural-analysis:palaeostress}
\emph{古应力}\index{古应力}是过去作用于岩石的应力状态，
只要能由其留下的构造推断出来即可。
实际上，它包括主应力的方向与
\emph{应力比}\index{应力比}
\[
  \Phi = \frac{\sigma_2 - \sigma_3}{\sigma_1 - \sigma_3}, \qquad 0 \le \Phi
  \le 1 .
\]
\emph{华莱士–博特假设}\index{华莱士–博特假设}指出，
每条断层沿均匀应力在其面上分解出的切应力
方向滑动，各断层彼此独立。\emph{断层滑动
反演}\index{断层滑动反演}是寻找主方向与应力比，
使其预测的滑动方向与一组
实测断层及擦痕拟合得最好。
\end{definition}

\begin{theorem}[平面上的滑动方向]\label{thm:b3:structural-analysis:wallace-bott}
设 $S$ 为应力张量，压应力取正，$\vec n$ 为
断层的向上单位法向量。记
\[
  \vec\tau = S\vec n - (\vec n\cdot S\vec n)\,\vec n .
\]
在\omterm{def:b3:structural-analysis:palaeostress}{华莱士–博特假设}下，上盘相对下盘
沿 $-\vec\tau$ 滑动，该力即上盘施加于
下盘的切向面力。此外：
\begin{enumerate}[label=(\alph*)]
  \item $\vec\tau$ 的方向在 $S$ 换为 $aS +
        bI$（$a > 0$）时不变；因此它只依赖于主方向与 $\Phi$，
        且约化张量 $S' =
        \mathrm{diag}(1, \Phi, 0)$ 在主轴系 $(\vec e_1, \vec
        e_2, \vec e_3)$ 中预测的滑动与 $S$ 相同；
  \item 若 $\vec n$ 垂直于某一主方向，则滑动
        亦垂直于该方向；
  \item 记 $n_k = \vec n\cdot\vec e_k$，所求方向即 $\vec\tau'
        = \vec\tau_0 + \Phi\,\vec\tau_1$ 的方向，关于 $\Phi$ 呈线性，其中 $\vec\tau_0
        = n_1(\vec e_1 - n_1\vec n)$，$\vec\tau_1 = n_2(\vec e_2 - n_2\vec
        n)$。
\end{enumerate}
\end{theorem}

\begin{proof}
压应力取正，上盘穿过断面向下盘施压，
断面的法向 $\vec n$ 指向上盘，上盘施加的
牵引力为 $-S\vec n$；其沿断面的分量即为 $-\vec\tau$。断层滑动时，
上盘运动方向即拖动下盘的方向，沿
$-\vec\tau$。（验证：取 $\sigma_1$ 竖直，$S = \sigma_1\vec e_z\vec
e_z^{\mathsf T}$，且 $-\vec\tau = -\sigma_1\cos\delta\,(\vec e_z -
\cos\delta\,\vec n)$ 沿倾向向下，为正断层。）
(a) $(aS + bI)\vec n - (\vec n\cdot(aS + bI)\vec n)\vec n = a\vec\tau + b(\vec
n - \vec n) = a\vec\tau$。令 $\sigma_k = \sigma_3 + (\sigma_1 -
\sigma_3)s_k$，取 $s = (1, \Phi, 0)$，则 $S = \sigma_3 I + (\sigma_1 -
\sigma_3)S'$。
(b) $S\vec n = \sum_k\sigma_kn_k\vec e_k$；若 $n_j = 0$，则 $S\vec n$
与 $\vec n$ 在 $\vec e_j$ 方向上均无分量，$\vec\tau$ 亦然。
(c) $S'\vec n = n_1\vec e_1 + \Phi n_2\vec e_2$，且 $\vec n\cdot S'\vec n =
n_1^2 + \Phi n_2^2$；将各项按是否含 $\Phi$ 合并即可。
\end{proof}

情形 (b) 即安德森的情形：包含 $\sigma_2$ 的断层
垂直于它滑动，与 $\Phi$ 无关，因而对\omterm{def:b3:structural-analysis:palaeostress}{应力
比}不提供信息。情形 (c) 则将实测滑动化为 $\Phi$。设 $\vec u$ 为
上盘的单位滑动矢量，由侧伏角给出，$\vec w = \vec n\times\vec
u$ 为断面内与之垂直的方向。预测滑动
平行于 $\vec u$，当 $(\vec\tau_0 + \Phi\vec\tau_1)\cdot\vec w = 0$ 时，
且由于 $\vec n\cdot\vec w = 0$，
\[
  \Phi = -\frac{n_1\,(\vec e_1\cdot\vec w)}{n_2\,(\vec e_2\cdot\vec w)} ,
\]
但须预测滑动随后指向 $+\vec u$ 而非 $-\vec u$。

当主方向未知时，每条断层只给出一个方程，
即其滑动在断面内的角度，未知量共四个：确定主轴方位的三个角，
以及 $\Phi$。产状不同的断层至少需要四条，
通常取几十条。反演使实测与预测之间
滑动角度的平方和最小，以最小二乘法
在网格上搜索，或如\cref{sec:b3:seismology-imaging:inverse}那样迭代求解。
不能拟合的断层构成第二个群体，或形成于另一时代。

下图给出瓦尔登丘陵的五条断层：主正断层
及另外四条较小的断层，均为正滑动。对 $\Phi$ 与
$\sigma_3$ 的倾伏向进行网格搜索（$\sigma_1$ 竖直），五条断层均能拟合到
四分之一度以内，$\sigma_3$ 水平，指向 \ang{120}。

\begin{omfigure}
\begin{tikzpicture}
\begin{axis}[width=0.56\linewidth, height=0.56\linewidth, axis equal image,
  hide axis, xmin=-1.15, xmax=1.15, ymin=-1.15, ymax=1.15, unbounded coords=jump]
% computed in figdata/bachelor-3/stereonet.py (Varden faults invented)
\addplot[gray!35, very thin] table[x=x, y=y] {figdata/out/bachelor-3-stereonet-schmidt.dat};
\addplot[black, thick, domain=0:360, samples=181] ({cos(x)}, {sin(x)});
\addplot[omDef, thick] table[x=x, y=y] {figdata/out/bachelor-3-stereonet-faults.dat};
\addplot[only marks, mark=*, mark size=1.8pt, omDef] table[x=x, y=y] {figdata/out/bachelor-3-stereonet-slips.dat};
\addplot[only marks, mark=square*, mark size=2.6pt, omThm] table[x=x, y=y] {figdata/out/bachelor-3-stereonet-s1.dat};
\addplot[only marks, mark=diamond*, mark size=3pt, omProp] table[x=x, y=y] {figdata/out/bachelor-3-stereonet-s2.dat};
\addplot[only marks, mark=triangle*, mark size=3pt, omMeth] table[x=x, y=y] {figdata/out/bachelor-3-stereonet-s3.dat};
\node[font=\scriptsize, above] at (axis cs:0,1.02) {北};
\node[omThm, font=\scriptsize, above left] at (axis cs:0,0) {$\sigma_1$};
\node[omProp, font=\scriptsize, above right] at (axis cs:0.5,0.87) {$\sigma_2$};
\node[omMeth, font=\scriptsize, below right] at (axis cs:0.87,-0.5) {$\sigma_3$};
\end{axis}
\end{tikzpicture}

\par\smallskip{\small 瓦尔登丘陵的断层滑动数据，等面积
网（下半球投影）：五个断面（大圆）及各上盘的滑动方向
（圆点），以及拟合最佳的
应力主值：$\sigma_1$ 竖直（方块），$\sigma_2$ 与 $\sigma_3$
水平（菱形与三角形，位于各轴的两端）。主
断层倾角为 \ang{60}，倾向 \ang{120}，包含 $\sigma_2$，
沿倾向正下方滑动；其余断层则
呈斜向滑动。}\label{fig:b3:structural-analysis:faults}
\end{omfigure}

本科一年级分册中的 P 轴和 T 轴并非主应力。它们
仅由断层面与滑动方向构造，位于包含滑动方向与法线的平面内，
与断层面成 \ang{45}。新断层位于内摩擦角约为 \ang{30} 的岩石中时，
与 $\sigma_1$ 约成 \ang{30}，
已活化的断层则可取任意角度。单个机制只能给出 P 轴和 T 轴；
只有将产状各异的许多断层联合反演，才能得到应力。这一假设本身有局限：
应力可能因断层而异，块体可能旋转，
相交的断层彼此阻碍滑动。

\section{剖面恢复}\label{sec:b3:structural-analysis:restoration}

剖面是一种解释：在地表与钻孔之间，
地质学家绘出的是可能的情形。本科二年级分册在讨论阿尔卑斯山的一章中
称某一剖面为平衡的，指其可恢复到未变形状态，
且地层长度与单元面积保持守恒，褶皱与逆冲断层扎根于滑脱面之上。
恢复本身即是检验。

\begin{definition}[剖面恢复、固定线、膝折带法]\label{def:b3:structural-analysis:restoration}
\emph{剖面恢复}\index{剖面恢复}在沿构造运动方向绘制的剖面上，
按事件发生的逆序，逐一消去其所示的变形，
并在剖面平面内保持
地层长度与单元面积不变。
\emph{固定线}\index{固定线}是横穿地层的一条线，位于地层彼此未发生滑动之处，
即未变形的前陆中，或对称褶皱的轴面上，
在恢复后的剖面中取其垂直于地层，长度均从此线量起。在
\emph{膝折带法}\index{膝折带法}中，褶皱画成等厚的直线翼，
两翼在尖锐的转折端相连，
其轴面平分两翼之间的夹角；
于是每层都是一条折线，
其长度为各直线段之和。
\end{definition}

一层岩层折成直线翼，两翼向两侧倾斜，倾角均为 $\delta$，
水平距离为 $L$，其长度为 $L/\cos\delta$，缩短量为
$S = L(1/\cos\delta - 1)$，长度的相对变化为 $\cos\delta - 1$。
面积给出平衡的另一半。

\begin{proposition}[至滑脱面的深度]\label{prop:b3:structural-analysis:detachment}
设某层位于水平滑脱面之上高度 $h$ 处，在剖面平面内缩短量为 $S$，
且该层与滑脱面之间的岩石保持面积不变，
剖面两端无物质穿越。则
褶皱层高出其原始高度的面积 $A$ 满足
\[
  A = S\,h , \qquad\text{ 从而 }\qquad h = \frac AS .
\]
\end{proposition}

\begin{proof}
变形之前，该层之下的岩石占据面积为 $L_0 h$ 的区域，
其长为 $L_0$。变形之后，同样的岩石位于长 $L_0 - S$ 的区域之上，
其上为该层，该层位于原始高度加隆升量处，其面积为
$(L_0 - S)h + A$。令两面积相等，得 $A = Sh$。
\end{proof}

两层岩层给出两个未知量：二者多出的面积之差，等于 $S$ 乘以两者的
竖直间距。下图给出了作图方法。

\begin{omfigure}
\begin{tikzpicture}[scale=1.0, >=stealth]
% (a) kink-band anticline over a detachment, drawn to balance: limbs 50 deg
% (tan 1.19175, 1/cos 1.55572), half-wavelength 2 -> rise 2.3835; S = 4/cos50 - 4
% = 2.2229, A = 4.767, h = A/S = 2.1445; competent layers 0.3 thick (0.4667
% vertically on the limbs); weak layer restored 1.5445 thick (areas 9.611 both)
\begin{scope}
\fill[omLithEvap] (0,0) -- (4,0) -- (4,1.2111) -- (2,3.5946) -- (0,1.2111) -- cycle;
\fill[omLithSand] (0,1.2111) -- (2,3.5946) -- (4,1.2111) -- (4,1.6778) -- (2,4.0613) -- (0,1.6778) -- cycle;
\fill[omLithLime] (0,1.6778) -- (2,4.0613) -- (4,1.6778) -- (4,2.1445) -- (2,4.528) -- (0,2.1445) -- cycle;
\draw (0,1.2111) -- (2,3.5946) -- (4,1.2111);
\draw (0,1.6778) -- (2,4.0613) -- (4,1.6778);
\draw[thick] (0,2.1445) -- (2,4.528) -- (4,2.1445);
\fill[pattern=north east lines, pattern color=omThm] (0,2.1445) -- (2,4.528) -- (4,2.1445) -- cycle;
\draw[omThm, dashed] (0,2.1445) -- (4,2.1445);
\node[omThm, font=\small, fill=white, inner sep=1.5pt] at (1.25,2.5) {$A$};
\draw[very thick] (-0.3,0) -- (4.3,0);
\draw[dashed, gray] (2,0) -- (2,4.8) node[above, font=\scriptsize, black] {轴面};
\draw[omMeth, very thick] (0,0) -- (0,2.1445) node[above left, font=\scriptsize] {固定线};
\draw[omMeth, very thick] (4,0) -- (4,2.1445) node[above right, font=\scriptsize] {固定线};
\draw[<->] (-0.35,0) -- (-0.35,2.1445) node[midway, left, font=\scriptsize] {$h$};
\draw[<->] (0,-0.35) -- (4,-0.35) node[midway, below, font=\scriptsize] {$L_0 - S$};
\node[font=\small] at (2,-1.0) {(a) 变形后};
\end{scope}
% (b) restored: length 4/cos 50 = 6.2229
\begin{scope}[xshift=5.6cm]
\fill[omLithEvap] (0,0) rectangle (6.2229,1.5445);
\fill[omLithSand] (0,1.5445) rectangle (6.2229,1.8445);
\fill[omLithLime] (0,1.8445) rectangle (6.2229,2.1445);
\draw (0,1.5445) -- (6.2229,1.5445);
\draw (0,1.8445) -- (6.2229,1.8445);
\draw[thick] (0,2.1445) -- (6.2229,2.1445);
\draw[very thick] (-0.3,0) -- (6.5,0);
\node[font=\scriptsize, above left] at (6.2229,0) {滑脱面};
\draw[omMeth, very thick] (0,0) -- (0,2.1445) node[above, font=\scriptsize] {固定线};
\draw[omMeth, very thick] (6.2229,0) -- (6.2229,2.1445) node[above, font=\scriptsize] {固定线};
\draw[dashed] (4,0) -- (4,2.9);
\draw[<->] (4,2.75) -- (6.2229,2.75) node[midway, above, font=\scriptsize] {$S$};
\draw[<->] (0,-0.35) -- (6.2229,-0.35) node[midway, below, font=\scriptsize] {$L_0$};
\node[font=\small] at (3.1,-1.0) {(b) 恢复后};
\end{scope}
\end{tikzpicture}

\par\smallskip{\small 膝折带褶皱的恢复，按平衡原则作图，翼部取
\ang{50}。(a) 两层刚性岩层折成等厚的直线翼，
覆于在核部增厚的软弱层之上，两条\omterm{def:b3:structural-analysis:restoration}{固定线}位于
两个向斜的轴面上，二者保持在原始
高度。(b) 恢复后的剖面：每层拉直，其长度沿翼部量取为
$L_0$，比变形后的剖面长出缩短量 $S$。
阴影面积 $A$ 为顶层高出其
原始高度的部分，它等于 $S$ 乘以其离滑脱面的高度 $h$；
软弱层在两幅图中
具有相同的面积。}\label{fig:b3:structural-analysis:restoration}
\end{omfigure}

穿过断层的剖面按事件的逆序恢复：
先恢复最年轻的构造。垂直断距为 $t_h$、倾角为 $\delta$ 的正断层，
在平行于其倾向的剖面上，水平断距为 $t_h/\tan\delta$，滑距为 $t_h/\sin\delta$，
上盘沿断层向上回移一个滑距。
若剖面与滑动方向斜交，则断层的倾角与位移
都是视值，一个看似平衡的剖面未必真的平衡：
这正是剖面须平行于运移方向作图的原因。

\begin{inthefield}[测量产状]
\begin{minipage}[t]{0.55\linewidth}\vspace{0pt}
测量层面时，地质学家将罗盘倾角计打开的盖子
平放于所选的光滑且具有代表性的表面上，
转动仪器本体，直至水准泡示平；此时指针给出
倾向。再将仪器侧立于最陡线方向，读出
倾角。读数记为倾向与倾角，例如 $075/13$，
并注明地点与表面的类型。用磁针读出的方位角
需以当地磁偏角加以校正。装有地质应用程序的
智能手机可一步测得同样的数据：将其平贴于
层面，加速度计给出重力方向，由此得出倾角，
磁力计则给出方位角。野外地质学家常将其与
指南针比对，因为附近的铁器和手机自身的电路
会干扰磁力计，因此他们对每个读数重复测量，
并用矢量平均法对法线取平均。
\end{minipage}\hfill
\begin{minipage}[t]{0.42\linewidth}\vspace{0pt}
\includegraphics[width=\linewidth]{images/book5/ai/b3-compass-clinometer.jpg}

{\scriptsize 罗盘倾角计靠放于倾斜的岩石面上，旁边
放着野外记录簿（AI 生成示意图）。}
\end{minipage}
\end{inthefield}

\section{练习}

\begin{exercise}[$\star$]\label{exo:b3:structural-analysis:1}
求产状为 $040/30$（倾向，
倾角）平面的极的倾伏向与倾伏角。求倾伏向为 \ang{315}、
倾伏角为 \ang{20} 的平面的产状。
\end{exercise}

\begin{exercise}[$\star$]\label{exo:b3:structural-analysis:2}
一条直线的倾伏角为 \ang{30}，它在\omterm{def:b3:structural-analysis:stereo}{等角网}与等面积网上
距圆心各多远？两网的半径均为 \qty{10}{cm}。
\end{exercise}

\begin{exercise}[$\star$]\label{exo:b3:structural-analysis:3}
某褶皱两翼的产状分别为 $090/30$ 和 $270/30$。计算
两翼极之间的夹角（用法线的点积求得）、翼间角，
以及该褶皱的类型。
\end{exercise}

\begin{exercise}[$\star$]\label{exo:b3:structural-analysis:4}
一断层产状为 $090/60$，其擦痕的侧伏角为 \ang{-45}。取走向矢量
$\vec s = (0, 1, 0)$，平面的下倾矢量为 $\vec d$，则
上盘的滑动矢量为 $\vec u = \cos r\,\vec s - \sin r\,\vec d$。计算
$\vec d$、$\vec u$，以及滑动矢量的倾伏向与倾伏角，并命名该断层。
\end{exercise}

\begin{exercise}[$\star\star$]\label{exo:b3:structural-analysis:5}
某褶皱两翼的产状分别为 $060/40$ 和 $250/30$。计算
褶皱轴的倾伏向与倾伏角，以及其翼间角。
\end{exercise}

\begin{exercise}[$\star\star$]\label{exo:b3:structural-analysis:6}
计算网中位于倾伏角大于 \ang{60} 的直线所对应圆内的面积占比，
分别在等面积网与\omterm{def:b3:structural-analysis:stereo}{等角网}上计算；
并计算下半球中倾伏角大于 \ang{60} 的方向所占的比例。
统计极点应使用哪种网？
\end{exercise}

\begin{exercise}[$\star\star$]\label{exo:b3:structural-analysis:7}
三个极位于竖直的东西向平面内，与竖直方向的夹角分别为 \ang{-20}、$0$ 和
\ang{20}，即 $\vec p = (\sin a, 0, -\cos a)$。计算它们的
\omterm{def:b3:structural-analysis:orientation-tensor}{方位张量}、本征值与特征向量，以及 $\beta$ 轴。为何
最小本征值恰为零？
\end{exercise}

\begin{exercise}[$\star\star$]\label{exo:b3:structural-analysis:8}
将椭球 $X : Y : Z = 3 : 1.5 : 1$ 标在 弗林图上，再将
$X = 1.8$、$Y = 1.1$ 的等体积椭球也标出。计算两者的 $k$，
并说明应变的类型。
\end{exercise}

\begin{exercise}[$\star\star$]\label{exo:b3:structural-analysis:9}
在平行于伸长方向切割的剖面中，原为圆形的鲕粒
变为轴比为 $2.25$ 的椭圆，且全部平行排列，每个椭圆的面积
保持不变。计算主伸长比，以及沿长轴
的伸长量。变形前，鲕粒中心之间的距离从未小于
\qty{0.4}{mm}，求 Fry 图中空白区域的半轴。
\end{exercise}

\begin{exercise}[$\star\star\star$]\label{exo:b3:structural-analysis:10}
在距天底 $\theta = \ang{60}$ 处，计算沿半径方向的伸长比
$\mathrm dr/(R\,\mathrm d\theta)$，以及沿圆周方向的伸长比 $r/(R\sin\theta)$，
两种投影均需计算。验证投影的定理，并给出
等面积网将该处的极小圆变成的椭圆的
轴比。
\end{exercise}

\begin{exercise}[$\star\star\star$]\label{exo:b3:structural-analysis:11}
在走滑应力体制中，$\sigma_1$ 水平指向 \ang{000}，
$\sigma_3$ 水平指向 \ang{090}，$\sigma_2$ 竖直。一竖直
断层走向为 \ang{020}。不计算 $\vec\tau$，证明其滑动
对任意 $\Phi$ 均为水平。然后计算 $-\vec\tau$（取 $\Phi = 0.4$），
平面的法线取 $\vec n = (\cos\ang{20}, -\sin\ang{20}, 0)$，并给出
滑动的方向。
\end{exercise}

\begin{exercise}[$\star\star\star$]\label{exo:b3:structural-analysis:12}
一条长 \qty{20}{km} 的剖面穿过一组滑脱褶皱。两个层位
相距 \qty{1.0}{km}，二者高出其原始高度的面积分别为
\qty{6.0}{km^2}（下层）和 \qty{7.5}{km^2}（上层）。计算缩短量、
滑脱面位于下层位之下的深度、剖面的原始长度，
以及剖面长度的相对变化量。
\end{exercise}

\section{问题：瓦尔登丘陵的褶皱与断层}

\begin{problem}[{周末问题——褶皱轴、断层组的应力与剖面恢复，
位于一个虚构地区}]
\label{pb:b3:structural-analysis:1}
% values: calc/ch08.py; figdata/bachelor-3/stereonet.py (data invented from region.py)
在瓦尔登丘陵（本系列的虚构地区）中，一套中生代盖层
覆于由片岩与花岗岩组成的基底之上。盖层被褶皱为直立的
背斜与向斜，两翼平直，随后被一条正断层切割。一位地质学家
在背斜的两翼共测得十二个层面的产状
（倾向/倾角，单位为度）：

\begin{center}
\small
\begin{tabular}{l|cccccc}
\hline
东翼 & 075/13 & 084/11 & 089/13 & 079/14 & 086/12 & 072/11 \\
西翼 & 283/11 & 272/13 & 281/14 & 288/10 & 276/13 & 274/12 \\
\hline
\end{tabular}
\end{center}

每翼六个读数的向上单位法向量之和分别为
$\vec S_E = (1.2583, 0.2017, 5.8604)$ 与 $\vec S_W = (-1.2445, 0.1886,
5.8636)$（东、北、上）。五条断层带有擦痕（倾向/倾角，
上盘滑动的侧伏角）：$120/60/{-90}$（主断层）、$150/55/{-76}$、
$090/65/{-108}$、$200/70/{-78}$ 和 $060/50/{-105}$。主断层的
垂直断距为 \qty{200}{m}。背斜与向斜的轴线相距 \qty{4}{km}，
因此褶皱自东向西每 \qty{8}{km} 重复一次。

\medskip\noindent\textbf{第一部分——极与 \omterm{def:b3:structural-analysis:pi}{β 轴}。}
\begin{enumerate}
  \item 给出读数 $075/13$ 的极的倾伏向与倾伏角，以及
        它到半径为
        \qty{10}{cm} 的等面积网中心的距离。
  \item 计算读数 $075/13$ 的向上单位法向量，并将其
        与 $\vec S_E$ 比较大小以作检验。
  \item 由 $\vec S_E$ 与 $\vec S_W$ 计算各翼的平均产状、
        平均合向量长度，以及各读数相对于其平均值的
        方均根离散度。
  \item 计算 $\vec S_E\times\vec S_W$，并求
        \omterm{def:b3:structural-analysis:pi}{β 轴}的倾伏向与倾伏角。
  \item 计算两平均法向量之间的夹角、翼间角，
        并判定褶皱的类别。
  \item 该对称褶皱的轴面包含 \omterm{def:b3:structural-analysis:pi}{β 轴}，且
        垂直于 $\vec m_E - \vec m_W$，即平均单位
        法向量之差。给出其产状；该褶皱是否直立？
  \item 给出 $\pi$ 环带的产状，并说明为何各极落在
        该环带上。
\end{enumerate}

\medskip\noindent\textbf{第二部分——\omterm{def:b3:structural-analysis:orientation-tensor}{方位张量}。}
十二个极的\omterm{def:b3:structural-analysis:orientation-tensor}{方位张量}为
\[
  T = \begin{pmatrix} 0.04400 & 0.00035 & 0.00109\\ 0.00035 & 0.00144 & 0.03179\\
  0.00109 & 0.03179 & 0.95456 \end{pmatrix},
\]
其本征值分别为 $0.95562$、$0.04400$ 和 $0.00038$。
\begin{enumerate}[resume]
  \item 验证 $T$ 的迹为 1。取 $\vec b = (0.0053, -0.9994,
        0.0333)$，即 $\vec S_E\times\vec S_W$ 方向的单位矢量，计算
        $\vec b^{\mathsf T}T\vec b$，该结果说明了 $\vec b$ 的什么？
  \item 计算各极与环带之间的方均根夹角。
  \item 由两翼的倾角解释 $\lambda_2$ 的数值。
  \item 在读数精度范围内，该褶皱是否为圆柱状？
  \item \omterm{def:b3:structural-analysis:pi}{β 轴}的倾伏角由环带的南–北向倾斜决定，
        各极对此倾斜的贡献取决于第 9 题求得的离散度。试估计
        倾伏角的标准误差。
\end{enumerate}

\medskip\noindent\textbf{第三部分——\omterm{def:b3:structural-analysis:palaeostress}{古应力}。}
\begin{enumerate}[resume]
  \item 主断层属于哪一类断层？按照安德森理论，给出
        $\sigma_1$、$\sigma_2$ 和 $\sigma_3$ 的方向。
  \item 利用关于滑动方向的定理证明，主断层
        沿其倾向向下滑动，对每一个 $\Phi$ 值都成立。这对
        $\Phi$ 说明了什么？
  \item 对断层 $150/55$，计算 $\vec n$ 及其分量 $n_1 =
        \vec n\cdot\vec e_1$ 与 $n_2 = \vec n\cdot\vec e_2$，其中 $\vec e_1$
        为竖直方向，$\vec e_2$ 为指向 \ang{030} 的水平方向。
  \item 其滑动给出 $\vec w = \vec n\times\vec u = (0.771, 0.605, 0.198)$。
        计算 $\Phi$。
  \item 其余三条小断层给出 $\Phi = 0.294$、$0.301$ 和
        $0.306$。求其平均值，并说明离散性的来源。
  \item $\Phi = 0.3$ 对 $\sigma_2$ 说明了什么？
  \item 计算主断层的 P 轴与 T 轴（它们与断面成 \ang{45}，
        位于其倾向面内），并将二者与 $\sigma_1$ 和
        $\sigma_3$ 比较。
\end{enumerate}

\medskip\noindent\textbf{第四部分——恢复东–西向剖面。}
\begin{enumerate}[resume]
  \item 计算东翼（$080.9/12.3$）在
        东–西向剖面上的视倾角。
  \item 取剖面上两翼的倾角为 \ang{12}，计算恢复后
        一个褶皱波长的长度、缩短量以及
        长度的相对变化量。
  \item 计算背斜顶部高出向斜的高度，以及
        一个波长内的多余面积，假设向斜保持在
        原始高度。
  \item 计算滑脱面位于褶皱层之下的深度。能否
        使褶皱滑脱于盖层底部，
        位于地下数百米处？
  \item 计算主断层在
        沿 \ang{120} 方向的剖面上的水平断距与滑距。
  \item 在东–西向剖面上，计算断层的视倾角与
        视水平断距。应恢复哪一个剖面？
  \item 断层与褶皱应按何种顺序恢复，且
        \omterm{def:b3:structural-analysis:restoration}{固定线}应放在何处？
  \item 由十二个读数恢复得到的瓦尔登褶皱轴的倾伏角
        是多少？
\end{enumerate}
\end{problem}
